\section{The Super Swap and Super Product Tests}\label{sec:super_swap}
In this section, we take another look at the product test as well as the swap test, considering one of the strongest possible generalization of the two.

We start with the more elementary swap test, which is a widely used to test if two quantum states, say $\ket{\psi}$ and $\ket{\phi}$, are equal.
One fundamental limitation of the swap test is that it always accepts with probability at least $1/2$ even if
the states are orthogonal. More precisely, its acceptance probability is $(1+\abs{\braket{\psi}{\phi}}^2)/2$.
Ideally, it would be much more useful to have a test with acceptance probability of $\abs{\braket{\psi}{\phi}}^2$, which is impossible with only one copy for each state.
In the presence of many unentangled copies of $\ket{\phi}$ but just a single copy of $\ket{\psi}$, we show that it
is possible to approach this goal with an arbitrarily small error overcoming the inherent limitation of the swap test. Therefore, we call this test the \emph{super swap} test
and we provide a description of it in~\cref{proc:super_swap}. In particular, this super swap test can be useful
when it is difficult to produce a state $\ket{\psi}$, but much easier to produce copies of $\ket{\phi}$
and we want the tester's acceptance probability to more accurately capture how close $\ket{\psi}$ is to $\ket{\phi}$.
In~\cref{sec:protocol}, the special state $\ket{\psi}$ will be a nonnegative amplitudes state which has a greater cost in the
context of complexity protocols there, whereas $\ket{\phi}$ will have general amplitudes being a cheaper resource in that context.  

\begin{algorithm}{SuperSwap}{($\ket{\psi},\ket{\phi}^{\otimes \ell}$)}
     \begin{enumerate}
           \item Project $\ket{\psi}\ket{\phi}^{\otimes \ell}$ onto the symmetric space $\vee^{\ell+1}(\C^{d})$.
           \item If the projection succeeds \emph{accept};        
            else \emph{reject}.
     \end{enumerate}
     \caption{Super Swap Test}\label{proc:super_swap}
\end{algorithm}

The acceptance probability of the super swap test is established next.

\begin{lemma}\label{lem:super_swap}
  The super swap test accepts with probability 
  \[\frac{\ell \cdot \abs{\braket{\psi}{\phi}}^2}{\ell+1} + \frac{1}{\ell+1}.\]
\end{lemma}

\begin{proof}
  Let $\Pi = (1/(\ell+1)!) \sum_{\pi \in \Sym_{\ell+1}} \pi$ be the projector onto $\vee^{\ell+1}(\C^{d})$. Indeed, we have
  \begin{align*}
    \bra{\psi}\bra{\phi}^{\otimes \ell} \Pi \ket{\psi}\ket{\phi}^{\otimes \ell} = \frac{1}{\ell+1} \braket{\psi}{\psi} \braket{\phi}{\phi}^{\ell} + \frac{\ell}{\ell+1} \abs{\braket{\psi}{\phi}}^2 \braket{\phi}{\phi}^{\ell-2}\mcom
  \end{align*}
  concluding the proof.
\end{proof}

At first glance, it may seem inconvenient to assume multiple ($\ell$-many) unentangled copies of $\ket{\phi}$.
However, due to our disentangler channel, we can enforce a distribution over product states $\ket{\phi}^{\otimes \ell}$ by assuming
only bipartite unentanglement.

Next we turn to the product test which checks whether a state is close to a $k$-partite product state~\cite{HM13}.
It has a similar drawback to the usual swap test, namely, it always accepts with probability at least $1/2$
even if the state $\ket{\psi}$ is very far from product. As before, we will arbitrarily improve the soundness of the product test
by having multiple unentangled copies. We call this new test the \emph{super product} test and we describe it in~\cref{proc:super_prod}.

\begin{algorithm}{SuperProduct}{($\ket{\psi},(\ket{\phi_1}\ldots\ket{\phi_k})^{\otimes \ell}$)}
     \begin{enumerate}
           \item Project $\ket{\psi}(\ket{\phi_1}\ldots\ket{\phi_k})^{\otimes \ell}$ onto the symmetric space $\vee^{\ell+1}((\C^{d})^{\otimes k})$.
          \item If the projection succeeds \emph{accept};        
            else \emph{reject}.
     \end{enumerate}
     \caption{Super Product Test}\label{proc:super_prod}
\end{algorithm}

\begin{lemma}
  The super product test accepts with probability
    \[ 
  \frac{\ell}{(\ell+1)} \cdot \abs{\braket{\psi}{\phi_1}\ldots\ket{\phi_k}}^2 + \frac{1}{(\ell+1)}\mper
    \] 
\end{lemma}

\begin{proof}
  We view each copy of the state $\ket{\phi_1}\ldots\ket{\phi_k}$ as a single state $\ket{\phi}$
  and apply the super swap test to $\ket{\psi}$ and $\ket{\phi}^{\otimes \ell}$. The acceptance  
  probability of the super product test now follows from~\cref{lem:super_swap}.
\end{proof}

Analogously, it may seem inconvenient to assume multiple ($\ell$-many) unentangled copies of $\ket{\phi_1} \ldots \ket{\phi_k}$.
However, that is not an issue by \cref{cor:separable-states}: We can enforce a distribution over product states $(\ket{\phi_1} \ldots \ket{\phi_k})^{\otimes \ell}$ by assuming only $2$ unentangled states. 
